\section{Coordinación temporal con TimerAction}

\subsection{Por qué hay que escalonar el arranque}

Gazebo necesita tiempo para inicializarse antes de que otro proceso pueda comunicarse con él. Si el nodo de spawn se lanza inmediatamente después de Gazebo, intentará conectarse a un simulador que aún está cargando y fallará. El mismo problema ocurre con el RSP: el spawn necesita leer el tópico robot\_description, que el RSP debe estar publicando. Y el nodo de control no debe arrancar hasta que el robot esté spawnado y con odometría disponible.

Esta cadena de dependencias temporales se resuelve con TimerAction: una acción que retrasa la ejecución de un conjunto de acciones hijas un número fijo de segundos desde que el launch file comenzó.

\subsection{Anatomía de TimerAction}

TimerAction recibe dos parámetros: period (el retardo en segundos, como float) y actions (una lista de acciones a ejecutar cuando el timer dispare). El Código~\ref{lst:timer_basico} muestra la forma básica:

\begin{figure*}[t]
\centering
\begin{lstlisting}[style=python,
    caption={Estructura básica de TimerAction.},
    label={lst:timer_basico}]
from launch.actions import TimerAction

actions.append(
    TimerAction(period=3.0, actions=[
        Node(
            package='gazebo_ros',
            executable='spawn_entity.py',
            ...
        )
    ])
)
\end{lstlisting}
\end{figure*}


El timer comienza a contar desde el momento en que el launch file empieza a ejecutarse, no desde que la acción anterior termina. Esto significa que todos los timers se lanzan en paralelo al inicio, y cada uno dispara de forma independiente cuando su period se cumple. No hay dependencia explícita entre timers.

\subsection{El esquema de retardos en consenso\_prom}

En consenso\_prom.launch.py se usa un esquema de retardos calculados para escalonar el spawn de $n$ robots sin sobrecargar Gazebo. El Código~\ref{lst:esquema_delays} muestra la lógica:

\begin{figure*}[t]
\centering
\begin{lstlisting}[style=python,
    caption={Esquema de retardos escalonados para $n$ robots.},
    label={lst:esquema_delays}]
for i, (x0, y0, yaw) in enumerate(poses):
    ns    = f'robot{i}'
    delay = 3.0 + i * 1.5
    # robot0: 3.0 s, robot1: 4.5 s, ..., robot7: 13.5 s

    actions.append(TimerAction(period=delay,
                               actions=[rsp_node]))
    actions.append(TimerAction(period=delay + 0.5,
                               actions=[spawn_node]))
    actions.append(TimerAction(period=delay + 1.5,
                               actions=[control_node]))
\end{lstlisting}
\end{figure*}


La secuencia para cada robot $i$ es: RSP en delay, spawn en delay + 0.5 s, nodo de control en delay + 1.5 s. Los 0.5 s entre RSP y spawn dan tiempo para que el tópico robot\_description esté disponible. El segundo extra antes del nodo de control da tiempo para que el robot aparezca en la simulación y empiece a publicar odometría.

Al separar robots con 1.5 s entre sí, los 8 robots completan su secuencia de arranque en unos 13.5 s, sin que ningún paso de spawn compita con otro por los recursos de Gazebo.

\subsection{Los nodos tardíos}

Los nodos que necesitan que todos los robots estén listos (el grafo de conectividad AIRE\footnote{AIRE es el nodo que calcula la conectividad entre robots basada en distancia euclidiana. Necesita que todos los robots estén publicando odometría antes de arrancar, de lo contrario podría reportar grafos incompletos durante los primeros ciclos.} y los plotters) se lanzan con un retardo calculado a partir del último robot:

\begin{figure*}[t]
\centering
\begin{lstlisting}[style=python,
    caption={Retardo para nodos que dependen de todos los robots.},
    label={lst:late_delay}]
late_delay = 3.0 + n * 1.5 + 2.0
# Para n=8: 3.0 + 12.0 + 2.0 = 17.0 s

actions.append(TimerAction(period=late_delay, actions=[
    Node(package='simulador', executable='aire', ...)
]))
\end{lstlisting}
\end{figure*}


El margen de 2.0 s adicional después del último robot garantiza que los nodos tardíos no arranquen antes de que el último robot de control esté operativo.
